\section{Ablations}
\label{sec:ablations}

We isolate the contribution of each acceleration mechanism on full UVG
720p with the Wan~2.1-T2V-1.3B backbone, $T\!=\!20$, $M\!=\!64$,
$K\!=\!16384$, $g_\text{scale}\!=\!3.0$ unless otherwise noted. Decode
time is reported per GOP.

\subsection{The target-domain mismatch and its fix}
\label{sec:abl-lpraw}

Table~\ref{tab:abl-mr} reports our main multi-resolution ablation on full
UVG ($7$~sequences~$\times$~$3$~GOPs) at 720p with Wan~2.1-T2V-1.3B.
All rows share the same codebook hyperparameters ($M\!=\!64$,
$K\!=\!16384$, $T_\text{sde}\!=\!17$) and the same single-resolution
baseline.

\begin{table}[t]
\centering
\caption{Multi-resolution coding with naive vs.\ unified low-stage target.
Full UVG ($7$~seq~$\times$~$3$~GOPs) at 720p, Wan~2.1-T2V-1.3B,
$S\!=\!2$ schedule ($480\!\to\!720$). $\Delta$ is relative to the
single-resolution baseline of $28.72$~dB at $0.00483$~BPP.}
\label{tab:abl-mr}
\small
\begin{tabular}{@{}l c c c@{}}
\toprule
Target & PSNR (dB) & BPP & $\Delta$PSNR \\
\midrule
Single-resolution baseline   & 28.72 & 0.00483 & --- \\
\midrule
MR naive + random fill       & 27.67 & 0.00453 & $-1.05$ \\
MR naive + DLC1 transition   & 28.07 & 0.00481 & $-0.65$ \\
\rowcolor{wacvblue!10}
MR \textbf{unified} + DLC1 (ours) & \textbf{28.50} & 0.00481 & \textbf{$-0.22$} \\
\bottomrule
\end{tabular}
\end{table}

\paragraph{Is the residual gap a codebook or a target problem?}
The naive scheme with random high-band fill loses $1.05$~dB. The DLC1
transition codebook recovers $\approx 0.4$~dB at $<1\%$ bitrate overhead,
but a residual $0.65$~dB gap remains. To check whether this gap is a
codebook-expressiveness problem or a codec-target problem, we ran a
diagnostic joint-oracle study on a two-sequence subset
(Bosphorus + HoneyBee, $3$ GOPs each), in which the encoder is allowed
to inject the ground-truth high-band content at the transition step
(non-deployable, zero bits transmitted for the oracle row).
Table~\ref{tab:abl-oracle} reports the result. \emph{The numbers in
Table~\ref{tab:abl-oracle} are not directly comparable to those of
Table~\ref{tab:abl-mr}: the diagnostic subset's single-resolution
baseline is $30.15$~dB rather than $28.72$~dB.}

\begin{table}[t]
\centering
\caption{Joint-oracle diagnostic on a small subset
(Bosphorus + HoneyBee, $3$ GOPs each), Wan~2.1-T2V-1.3B at 720p,
$S\!=\!2$. The subset's single-resolution baseline is $30.15$~dB.
Numbers are not directly comparable to Table~\ref{tab:abl-mr}.}
\label{tab:abl-oracle}
\small
\begin{tabular}{@{}l c c@{}}
\toprule
Target & PSNR (dB) & $\Delta$ vs $30.15$ \\
\midrule
Single-resolution (subset baseline) & 30.15 & --- \\
MR naive + random        & 28.46 & $-1.69$ \\
MR naive + DLC1          & 28.72 & $-1.43$ \\
MR naive + joint oracle  & 28.89 & $-1.26$ \\[2pt]
\rowcolor{wacvblue!10}
MR \textbf{unified} + DLC1 (ours) & 30.05 & $-0.10$ \\
\bottomrule
\end{tabular}
\end{table}

Under the naive target, even the joint oracle is capped at $-1.26$~dB on
the diagnostic subset: the gap is intrinsic to the codec target the
encoder is matching against, not to codebook expressiveness. Under the
unified target, the same diagnostic gap collapses to $-0.10$~dB. Both
observations are consistent with the full-UVG numbers in
Table~\ref{tab:abl-mr}.

\paragraph{Magnitude alignment is not the fix.}
We tried two natural variants of the unified target that rescale
$\mathcal{T}_L(\mathrm{VAE}(\mathbf{V}))$ to match the per-channel std
(or norm) of the naive target before use, intending to ``soften'' the
domain gap. Both \emph{hurt}: by $\approx 0.5$~dB on the diagnostic
subset relative to the raw unified target. Codebook-matching is
spectrally sensitive, and rescaling re-introduces the domain mismatch
in a more subtle form.

\paragraph{The fix is essentially one line of code.}
Replacing Eq.~\ref{eq:naive-target} with Eq.~\ref{eq:lp-target} collapses
the full-UVG gap from $-0.65$ (naive + DLC1) to $-0.22$~dB at the same
bitrate. The diagnostic-subset joint-oracle ceiling under the unified
target is $-0.10$~dB, showing that once the target is corrected, transition
high-band coding is no longer the dominant bottleneck.

\subsection{Predictive skip: caching choice and schedule}
\label{sec:abl-skip}

Table~\ref{tab:abl-skip} compares no caching, velocity caching, and
clean-endpoint caching on UVG-720p with Wan~2.1-T2V-1.3B, at the fixed
alternating $p\!=\!2$ schedule (every other SDE step skipped) with
transition warmup after the multi-resolution lift.

\begin{table}[t]
\centering
\caption{Caching choice for the per-step generator query, $p\!=\!2$
alternating schedule, full UVG 720p, multi-resolution
$S\!=\!2$. Time is per-GOP decoder wall-clock; $\Delta$PSNR is relative
to no caching.}
\label{tab:abl-skip}
\small
\begin{tabular}{@{}l c c c@{}}
\toprule
Cache & PSNR (dB) & Time/GOP & $\Delta$PSNR \\
\midrule
None (full $T$ generator calls)    & 28.50 & 25s & --- \\
Velocity cache (Eq.~\ref{eq:u-cache})       & 26.71 & 19s & $-1.79$ \\
\rowcolor{wacvblue!10}
\textbf{Clean-endpoint cache} (ours, Eq.~\ref{eq:u-tilde}) & 28.45 & 19s & $-0.05$ \\
\bottomrule
\end{tabular}
\end{table}

The clean-endpoint cache reduces decoder time by $1.32\times$ at a
$0.05$~dB PSNR cost; the same skip rate with naive velocity caching costs
$1.79$~dB---a $36\times$ larger PSNR drop. The two schemes are otherwise
identical (same skip schedule, same number of generator calls); the
difference comes entirely from \emph{what} is cached. The velocity cache
is computed at a stale state and the trajectory has moved away from
$\mathbf{x}_{t_k}$ by both the deterministic drift and the codebook
innovation between steps $k$ and $k\!+\!1$. The endpoint cache is
geometrically invariant along the deterministic part of the trajectory,
so only the codebook-induced shift between steps remains as approximation
error.

\paragraph{The one-step radius.}
Table~\ref{tab:abl-cap} ablates the maximum allowed number of
consecutive skips, holding the average skip rate fixed via an adaptive
threshold (we transmit a 1-bit-per-SDE-step skip mask in the bitstream
when adaptivity is on).

\begin{table}[t]
\centering
\caption{Adaptive predictive skip with a per-stage cap on consecutive
skips. HoneyBee, single GOP, 720p, Wan-T2V-1.3B. Threshold tuned to
keep average skip rate near $50\%$.}
\label{tab:abl-cap}
\small
\begin{tabular}{@{}l c c c@{}}
\toprule
Schedule & PSNR (dB) & Time/GOP & Notes \\
\midrule
Fixed alternating ($p\!=\!2$, cap $\equiv 1$) & 29.47 & 18s & ours \\
Adaptive, cap $=\!1$, thr $0.30$               & 29.48 & 19s & robustness mode \\
Adaptive, cap $=\!2$, thr $0.15$ (triggered)   & 29.08 & 17s & $-0.39$ \\
Adaptive, no cap, thr $0.40$ (triggered)       & 25.95 & 13s & $-3.52$ \\
\bottomrule
\end{tabular}
\end{table}

Whenever the cap of $2$ is actually triggered (i.e., the adaptive
controller picks two consecutive skips), quality drops sharply. Removing
the cap entirely allows long skip runs and quality collapses by
$3.5$~dB. With a hard cap of $1$, the adaptive controller can only
\emph{cancel} a would-be-skip when the running per-step error is too
large, never schedule a faster pattern; quality therefore matches the
fixed alternating schedule but speed is bounded above by the same.
\emph{Adaptive caching with cap $\equiv 1$ is a robustness mechanism,
not a new acceleration frontier}. The fixed alternating schedule is
Pareto-optimal under the clean-endpoint caching scheme.

\subsection{Stage count and transition placement (diagnostic subset)}
\label{sec:abl-stages}

Table~\ref{tab:abl-stages} ablates the number of multi-resolution stages
$S$ and the placement of transition steps. The schedule sweep is
relatively expensive (six configurations~$\times$~multiple
sequences), so we run it on a \emph{diagnostic subset}---Bosphorus and
HoneyBee, $3$~GOPs each---rather than on full UVG. \textbf{The single-
resolution baseline for this subset is $30.15$~dB, distinct from the
$28.72$~dB full-UVG baseline used in Table~\ref{tab:abl-mr}}; the two
tables answer different questions on different data and their numbers
should not be combined. All entries below use the unified target
(Eq.~\ref{eq:lp-target}) and the DLC1 transition codebook.

\begin{table}[t]
\centering
\caption{\textbf{Diagnostic-subset} stage-count and transition-placement
sweep (Bosphorus + HoneyBee, $3$~GOPs each, Wan-T2V-1.3B at $720$p).
$\tau$ lists the SDE step indices at which each transition occurs.
$\Delta$PSNR is relative to this subset's single-resolution baseline of
$30.15$~dB (not the full-UVG $28.72$~dB baseline of
Table~\ref{tab:abl-mr}).}
\label{tab:abl-stages}
\small
\begin{tabular}{@{}l c c c c@{}}
\toprule
Schedule & $S$ & $\tau$ & PSNR (dB) & BPP \\
\midrule
Single resolution         & 1 & --- & 30.15 & 0.00483 \\
\rowcolor{wacvblue!10}
$480 \!\to\! 720$            & 2 & $\{8\}$    & 28.72 & 0.00481 \\
$480 \!\to\! 720$ (early $\tau$) & 2 & $\{6\}$ & 28.51 & 0.00481 \\
$480 \!\to\! 720$ (late $\tau$)  & 2 & $\{10\}$ & 28.65 & 0.00481 \\
$272 \!\to\! 480 \!\to\! 720$     & 3 & $\{4,12\}$ & 27.76 & 0.00506 \\
$272 \!\to\! 480 \!\to\! 720$ (oracle) & 3 & $\{4,12\}$ & 28.07 & 0.00451 \\
\bottomrule
\end{tabular}
\end{table}

\paragraph{$S\!=\!2$ is the sweet spot on Wan~2.1.}
Two stages with a transition at the midpoint dominate $S\!=\!3$ schedules
that include a $272$p floor. The drop at $S\!=\!3$ is not a coding
artifact: even the joint oracle at $S\!=\!3$ ($28.07$~dB) cannot reach
the $S\!=\!2$ codebook result ($28.72$~dB). The cause is that Wan~2.1
was not trained at latents below the $480$p band; running the first few
denoising steps on an even smaller latent operates the generator
out-of-distribution, and the spectrally lifted state at transition is no
longer a good initialization for the higher-resolution stage. Moving the
low floor up to $480$p eliminates the problem; further increasing $S$
within $[480, 720]$ adds transition overhead without quality benefit.

\paragraph{Transition timing.}
At $S\!=\!2$, sweeping the transition step in $\{6, 8, 10\}$ yields
$28.51 / 28.72 / 28.65$~dB respectively. The step-$8$ choice
(approximately midway through the SDE phase) is best by a small margin
($+0.21$ over step-$6$). This matches the spectral diffusion
intuition~\cite{xiao2026spd}: too early a switch deprives the low stage
of useful denoising, too late a switch wastes generator capacity at high
resolution.

\subsection{Vectorized codebook RNG: implementation, not algorithm}
\label{sec:abl-vec}

The Turbo-DDCM~\cite{vaisman2025turboddcm} codebook construction
generates $K$ atoms per step per latent frame by independent
\texttt{randn}\,$(D)$ calls. In our implementation, replacing these with
a single batched \texttt{randn}\,$(K, D)$ per step reduces decoder
wall-clock by $\approx 1.9\times$. The batched call produces a
\emph{different} per-atom RNG sequence than the sequential call, so the
two settings emit different codebook indices for the same input; what
they preserve is the bit budget and the deterministic encoder/decoder
reproducibility (both sides agree as long as both use the same
batching). This is an engineering fix---it does not
appear in the bitstream---but we report it for transparency: all timing
numbers in this paper, including the published GVCC baseline rows in
Table~\ref{tab:time}, include this fix unless explicitly marked
\textsc{vanilla}.
